\section{Introduction}
KV state grows with context while attention concentrates on a subset of
positions. KV-management methods exploit this through
accumulated attention~\cite{h2o}, recent-query observations~\cite{snapkv},
query--key page bounds~\cite{quest}, cross-layer information~\cite{infinigen} and
learned retention~\cite{locret}. Better prediction need not identify a better
retention policy: selection changes later model states and their outputs.
Traceable evaluation of inference state under a budget connects this study to
D2AI's infrastructure, provenance and benchmark scope.

We audit which signals predict future attention, whether predictor complexity
helps, and how proxy rankings relate to answer quality. Offline labels,
served-oracle diagnostics and closed-loop answer scores are distinct objects;
the runtime scorer also reconstructs its features.

\textbf{Contributions.} (1)~\emph{A protocol and benchmark.} KVSalienceBench fixes
future-attention labels, selector semantics (Table~\ref{tab:protocol}), request-
and source-prompt splits, tie-aware per-decision metrics and versioned collection.
(2)~\emph{Controlled findings.} Two attention-magnitude views
support a three-parameter logistic scorer within 0.006 AUC of LightGBM on v2 traces
without a calibration step. At the per-decision top-decile budget, the within-layer
signal recovers more future top-decile blocks than either; this reversal shrinks
at 20\% retention. A separate pinned bfloat16 masked-loop rerun finds
the reconstructed scorer and accumulator mean-F1 equivalent within a pre-specified
$\pm0.02$ margin over 14 model--dataset cells (TOST $p=0.0007$).
A bridge measures feature reconstruction (\S\ref{sec:realization}).
(3)~\emph{An audit and its repair.} We diagnose corrupted float16 Qwen2.5
generation, rerun with a pinned simulator, audit capacity effects from logged counts,
and report physical-KV parity levels. Code, tests and provenance accompany the paper.
